\footnotesize
\setlength{\tabcolsep}{4pt}
\begin{longtable}{p{3.1cm}p{4.4cm}>{\raggedleft\arraybackslash}p{1.9cm}>{\raggedleft\arraybackslash}p{1.9cm}>{\raggedleft\arraybackslash}p{1.9cm}}
\caption{\textbf{Table S29.} Effect of gradual convergence to 75\% under three appraisals: the dynamic model against the same-year status quo, a static calculation that applies the long-run elasticities at once to the 2025 or the 2060 status quo population, and a comparison against the status quo frozen at its 2025 value. Primary simulation run.}\\
\toprule
\textbf{Appraisal} & \textbf{Comparison} & \textbf{Consumption change, \%} & \textbf{Prevalence change, pp} & \textbf{Change in T1-T3 gap, pp} \\
\midrule
\endfirsthead
\toprule
\textbf{Appraisal} & \textbf{Comparison} & \textbf{Consumption change, \%} & \textbf{Prevalence change, pp} & \textbf{Change in T1-T3 gap, pp} \\
\midrule
\endhead
\bottomrule
\endfoot
Dynamic model (main text) & Policy 2060 vs status quo 2060 & -14.38 & -1.07 & -1.82 \\
Static calculation & Status quo 2025 at the long-run price, adjusted at once & -17.10 & -2.02 & -3.12 \\
Static calculation & Status quo 2060 at the long-run price, adjusted at once & -17.65 & -1.64 & -2.82 \\
Frozen baseline & Policy 2060 vs status quo 2025 & -39.16 & -5.97 & 0.34 \\
\end{longtable}
\normalsize
